\begin{abstract}
Parsing is sequential in its classical formulation, yet state-of-the-art JSON
parsers such as \textsf{simdjson} and \textsf{cuJSON} run an order of magnitude
faster than textbook parsers by exploiting a small set of data-parallel
primitives.  These primitives are usually presented as format-specific tricks
(\textsf{prefix\_xor}, depth-sort bracket matching, \ldots), which makes them
hard to reuse and impossible to derive.

We show that they are instances of a single algebra, the \emph{scan algebra} $\SA$, and that
a grammar can be compiled to a data-parallel parser by a \emph{mechanical
rewrite} of a naive branch program into branchless masked data flow:
a stateful byte recurrence becomes a prefix product in a \emph{transformation
monoid} (a parallel scan), a filter becomes a \emph{compaction}, and the parser
stack becomes a \emph{well-nested reduction}.  Masking is derived from a
compiled automaton, not from the syntactic shape of rules; a specialization is
emitted only when it is provably equivalent to the automaton.

We build a compiler that turns a grammar into a CUDA implementation, and
evaluate it on JSON, CSV, TOML, XML, INI and a C-like language, on two testbeds
(a laptop GPU and an A100 node).  Separating transport from compute, the
generated device-resident scanners match or exceed domain SOTA: on CSV they
reach $3.6\times$ \textsf{simdcsv} on the laptop and $23\times$ on the A100,
and on TOML/INI/XML and C-like---none of which has a SIMD or GPU parser---they
run far ahead of the best CPU parsers.  Adding data movement is decisive: the
host-boundary chunked path is PCIe-limited, staying below \textsf{simdcsv} on
CSV on the laptop while surpassing it on the A100 node, whose shared CPU is the
bottleneck.
\end{abstract}
